\documentclass[11pt]{article}

\usepackage[T1]{fontenc}
\usepackage[utf8]{inputenc}
\usepackage[margin=1in]{geometry}
\usepackage{microtype}
\usepackage{amsmath,amssymb}
\usepackage{booktabs}
\usepackage{tabularx}
\usepackage{array}
\usepackage{float}
\usepackage{enumitem}
\usepackage{xcolor}
\usepackage{hyperref}

\definecolor{linkblue}{RGB}{32,91,155}
\hypersetup{
  colorlinks=true,
  linkcolor=linkblue,
  citecolor=linkblue,
  urlcolor=linkblue,
  pdftitle={A Causal Audit Layer for Deletion in Co-LMLM}
}

\setlist[itemize]{leftmargin=1.4em,itemsep=0.2em,topsep=0.3em}
\setlist[enumerate]{leftmargin=1.6em,itemsep=0.2em,topsep=0.3em}
\newcolumntype{Y}{>{\raggedright\arraybackslash}X}

\newcommand{\FULL}{\textsc{Full}}
\newcommand{\DELON}{\textsc{Del-On}}
\newcommand{\DELOFF}{\textsc{Del-Off}}
\newcommand{\colmlm}{Co-LMLM}
\newcommand{\pct}[1]{#1\%}

\title{\textbf{Co-LMLM Deletion Audit}\\
\large What We Found Before Testing Identity Clustering}
\author{Arya Raeesi \and Hanna Roed}
\date{Working note, 25 July 2026}

\begin{document}
\maketitle

\section*{Short version}

We ran a three-state deletion audit on the public
\texttt{lil-lab/CoLMLM-360M-FW} model and its released Wikipedia index. We
used 1,500 examples from each of T-REx, PopQA, Google-RE, CounterFact, and
ZsRE. Our detailed T-REx cohort contains 499 cases where the unmodified model
is correct and retrieves an answer-mentioning span before the answer is
visible.

Deleting only that retrieved row leaves \pct{39.1} of the T-REx answers
intact. Deleting its cosine neighborhood leaves \pct{31.1}. A filter supplied
with the gold answer and aliases reduces survival to \pct{6.8}, and combining
that filter with geometry reduces it to \pct{5.8}. This last result is a useful
efficacy upper baseline, but it is not a deployable deletion rule: it uses the
ground-truth answer at inference time, can remove unrelated facts that share
the same answer, and guarantees that our answer-mention trace check is zero.

The same combined policy leaves between \pct{0.3} and \pct{5.8} intact across
the five datasets. On T-REx, \pct{23.2} of the cohort remains answerable when
factual retrieval is disabled, and retrieval changes a correct
retrieval-disabled answer into a wrong answer in \pct{18.4} of cases. A
synthetic answer-bearing row placed just outside the cosine boundary also
restores many answers. Our main conclusion is therefore narrower than
``semantic deletion works'': a database row, a source, a cosine neighborhood,
and an answer string are different from a proposition identity. This note is
the audit before identity clustering; identity clustering itself is not
implemented here.

\section{Scope and relation to the Co-LMLM paper}

\colmlm externalizes free-form factual spans into a continuous-key,
text-value index and retrieves them from hidden-state queries
\cite{feldman2026colmlm}. The paper reports aggregate forgetting on TOFU
\cite{maini2024tofu}. Our run complements that experiment; it does not
reproduce it. The TOFU experiment fine-tunes on synthetic biographies, adds
those biographies to the knowledge base, and removes entries for a designated
forget set. We instead audit naturally occurring retrieval paths in the
released Wikipedia index.

The question here is: after a deletion request, does the answer still arrive
through another memory entry, or can the model and prompt produce it without
factual retrieval? This separates failures of the deletion set from answers
that remain available through the prompt, model parameters, or general
language patterns.

This is also the intended step before identity-key research. We first test
several current meanings of ``delete this fact''---one row, one source, a
cosine neighborhood, or a matching answer string. We can then compare those
baselines with a persistent proposition identity.

\section{Audit protocol}

Our design adapts the causal audit used for relational LMLMs
\cite{raeesi2026audit}. For each fact, the prompt, model, decoding, and
parameters stay fixed while memory access changes:

\begin{center}
\begin{tabularx}{0.92\textwidth}{p{0.14\textwidth}Y}
\toprule
State & Intervention \\
\midrule
\FULL & Memory unchanged; retrieval enabled. \\
\DELON & Target deletion set hidden; retrieval enabled. \\
\DELOFF & Target deletion set hidden; factual retrieval disabled. \\
\bottomrule
\end{tabularx}
\end{center}

Let \(C_s(f)\) indicate whether the output in state \(s\) contains the
normalized gold answer or an accepted alias as a complete phrase. Conditional
on \FULL\ correctness, we report

\begin{align}
L(f) &= C_{\DELOFF}(f),
  &&\text{retrieval-disabled answerability},\\
R(f) &= C_{\DELON}(f)\bigl(1-C_{\DELOFF}(f)\bigr),
  &&\text{retrieval-mediated correctness},\\
I(f) &= \bigl(1-C_{\DELON}(f)\bigr)C_{\DELOFF}(f),
  &&\text{retrieval interference},\\
S(f) &= C_{\DELON}(f),
  &&\text{post-deletion survival}.
\end{align}

The correctness rule is directional: the complete normalized answer or alias
must appear in the model output. A shorter prediction contained inside the
answer does not count. For example, ``Black'' is not accepted for
``Blackpool.''

\subsection{Cohort construction}

We audit the first retrieval event selected by the backend that satisfies all
of the following: it returned an entry ID, the entry passes our normalized
answer-mention heuristic, and the answer was not already visible before that
retrieval event. The heuristic is not a proposition verifier. It checks
whether the retrieved text contains the normalized answer or alias; it also
accepts a retrieved span of at least three characters contained inside the
answer. In the T-REx cohort, 570 entries pass by direct answer mention and two
pass only by this reverse check.

\begin{table}[H]
\centering
\caption{Construction of the auditable cohorts. ``Pre-answer mention'' means
that the selected retrieved span passed the answer-mention heuristic before
the answer was visible. The primary cohort also requires a correct \FULL\
answer.}
\label{tab:coverage}
\small
\begin{tabular}{lrrrr}
\toprule
Dataset & Prompts & Selected entry & Pre-answer mention & Primary \\
\midrule
T-REx & 1,500 & 1,095 & 572 & 499 \\
PopQA & 1,500 & 1,438 & 618 & 468 \\
Google-RE & 1,500 & 1,475 & 787 & 759 \\
CounterFact & 1,500 & 1,102 & 439 & 358 \\
ZsRE & 1,500 & 1,042 & 535 & 337 \\
\bottomrule
\end{tabular}
\end{table}

For T-REx, the 928 exclusions are 405 prompts with no selected entry, 509
whose selected span does not pass the answer-mention heuristic, and 14 where
the answer is already visible before the answer-mentioning event. The
unmodified model is correct on 746 of all 1,500 T-REx prompts under our
scoring rule. The 499-case primary cohort is \pct{33.3} of the full slice.

\subsection{Run configuration}

All five runs use greedy decoding with at most 12 new tokens, retrieval
similarity threshold \(0.7\), the public
\texttt{lil-lab/CoLMLM-360M-FW} checkpoint, and the released approximate
Wikipedia index of about 240M entries. They do not use the 2.2B-entry
Wikipedia--FineWeb-Edu index \cite{feldman2026colmlm}. T-REx and PopQA use
stable seeded slices; Google-RE is balanced across place of birth and place
of death; CounterFact and ZsRE use the first 1,500 eligible records. The
T-REx sample, 12-token limit, and containment metric differ from the
paper's published benchmark protocol, so the \FULL\ score above is not a
replication of the paper's T-REx result.

Deletion is implemented through reversible search-time filtering rather than
physical index rewriting. Every successful retrieval event was aligned to its
generated \texttt{<FACT>} block, and the geometric closure was centered on the
same audited event used to select the deletion target.

\section{Results}

\subsection{Five-dataset summary}

The standard policy combines geometric deletion at \(\rho=0.85\) with the
gold-answer string filter. Table~\ref{tab:datasets} reports the paired outcomes
on each primary cohort.

\begin{table}[H]
\centering
\caption{Three-state results under geometric plus gold-answer
string filtering. All rates are conditional on \FULL\ correctness and
pre-answer answer-mention retrieval.}
\label{tab:datasets}
\small
\begin{tabular}{lrrrrr}
\toprule
Dataset & Primary & Survival \(S\) & Disabled \(L\) & Mediated \(R\) &
Interference \(I\) \\
\midrule
T-REx & 499 & \pct{5.8} & \pct{23.2} & \pct{1.0} & \pct{18.4} \\
PopQA & 468 & \pct{2.4} & \pct{8.1} & \pct{1.3} & \pct{7.1} \\
Google-RE & 759 & \pct{0.3} & \pct{0.3} & \pct{0.0} & \pct{0.0} \\
CounterFact & 358 & \pct{4.7} & \pct{20.4} & \pct{0.3} & \pct{15.9} \\
ZsRE & 337 & \pct{0.6} & \pct{0.9} & \pct{0.3} & \pct{0.6} \\
\bottomrule
\end{tabular}
\end{table}

The size of the retrieval-disabled channel varies strongly by prompt format:
it is substantial on T-REx and CounterFact, smaller on PopQA, and nearly
absent on Google-RE and ZsRE. The deletion-policy result is more consistent:
combined-policy survival is at most \pct{5.8} on every dataset. Because the
policy receives the gold answer, this consistency measures end-to-end
efficacy of an oracle string filter, not successful proposition resolution.

\subsection{T-REx channel decomposition}

\begin{table}[H]
\centering
\caption{Paired T-REx outcomes under geometric plus gold-answer string
filtering at \(\rho=0.85\), restricted to the 499-case primary cohort.}
\label{tab:three-state}
\begin{tabularx}{0.93\textwidth}{Yrr}
\toprule
Outcome & Count & Rate \\
\midrule
Wrong in both \DELON\ and \DELOFF & 378 & \pct{75.8} \\
Correct only in \DELON: retrieval-mediated \(R\) & 5 & \pct{1.0} \\
Correct only in \DELOFF: retrieval interference \(I\) & 92 & \pct{18.4} \\
Correct in both post-deletion states & 24 & \pct{4.8} \\
\midrule
Any \DELON\ correctness: survival \(S\) & 29 & \pct{5.8} \\
Any \DELOFF\ correctness: \(L\) & 116 & \pct{23.2} \\
\bottomrule
\end{tabularx}
\end{table}

The standard policy leaves 29 of 499 answers intact
(\pct{5.8}, Wilson 95\% interval \([4.1,8.2]\)). Only five are correct
uniquely with retrieval. By contrast, 116 are correct without factual
retrieval. In 92 cases, retrieval changes a correct retrieval-disabled answer
into a wrong answer.

All 29 \DELON-correct traces lack a retained answer mention. This is expected
by construction: the same answer-mention rule is used both to filter
candidates and to inspect the trace. It is therefore a sanity check on the
implementation, not independent evidence that all semantic support was
removed.

The \DELOFF\ result is similar under the stricter control that forbids the
retrieval token: 110 of 499 are correct (\pct{22.0}), compared with 116
(\pct{23.2}) under null retrieval. The paired difference is 1.2 percentage
points (\(p=0.362\)). We call this retrieval-disabled answerability rather
than pure weight memorization because the prompt and general language
patterns remain available.

\subsection{Deletion policy and observed scope}

We compare six end-to-end policies. Selected-entry deletion hides the audited
row. Provenance deletion also hides entries from its source. Geometric
deletion hides keys with cosine similarity at least \(0.85\). Gold-answer
string filtering receives the target answer and aliases at inference time and
hides matching candidates. The combined and hybrid policies take unions of
these rules.

\begin{table}[H]
\centering
\caption{T-REx policy comparison on the 499-case primary cohort. Scope
statistics use all 572 pre-answer answer-mention targets. ``Observed excluded
IDs'' is the union of the saved manifest and candidate IDs encountered and
filtered during generation; it is a lower bound on total policy scope.
Provenance scope is source-wide and cannot be counted from the saved trace.
\(\dagger\) Zero answer mentions are guaranteed by the gold-answer filter.}
\label{tab:policies}
\scriptsize
\begin{tabular}{lrrrrrr}
\toprule
Policy & Survival & Efficacy & Mention in survivor & Mean & Median & Max. \\
\midrule
Selected entry & \pct{39.1} & \pct{60.9} & 178/195 & 1.0 & 1 & 1 \\
Provenance & \pct{39.7} & \pct{60.3} & 185/198 & --- & --- & --- \\
Geometric, \(\rho=0.85\) & \pct{31.1} & \pct{68.9} & 139/155 & 2.6 & 1 & 183 \\
Gold-answer string filter & \pct{6.8} & \pct{93.2} & \(0/34^\dagger\) &
79.8 & 23 & 2,120 \\
Geometric \(+\) gold-answer & \pct{5.8} & \pct{94.2} &
\(0/29^\dagger\) & 86.5 & 24 & 2,835 \\
Hybrid & \pct{5.8} & \pct{94.2} & \(0/29^\dagger\) &
88.4 & 25 & 2,835 \\
\bottomrule
\end{tabular}
\end{table}

Deleting one observed row is not the same as deleting the fact: 178 of the
195 selected-entry survivors retrieve another answer-mentioning span.
Geometric closure removes more, but 139 of its 155 survivors still retrieve
an answer mention. The same record--fact mismatch has been identified as a
structural problem for selective forgetting in RAG systems
\cite{tavakoli2026deletion}.

The gold-answer filter produces the large efficacy gain. There are 124 cases
that survive geometric deletion but not gold-answer filtering, versus three
in the reverse direction (\(p<5\times10^{-33}\)). This is an oracle upper
baseline, not a semantic policy. Different subjects and relations can share
the same answer, so a string rule can remove lawful and unrelated facts.

The saved top-500 value manifests contain a mean of 60.0 entries per target
(median 22, maximum 494). During generation, the filter also applies to newly
encountered candidates. Taking the union of saved and dynamically observed
IDs raises the mean to 79.8, the median to 23, and the maximum to 2,120.
These are still observed lower bounds, not the number of matching entries in
the whole 240M-entry index.

There is also an approximate-search confound. Selected-entry and geometric
policies search at shallow, target-dependent depths, while value and
provenance filtering search about 4,097 candidates at a time. On a compressed
approximate index, changing search depth can change the apparent top result.
The table therefore compares complete inference policies; it does not isolate
deletion-set membership alone. A fixed-depth validation subset is needed for
that stronger causal comparison.

\subsection{Geometric boundary stress test}

For each T-REx target, this test inserts one synthetic answer-bearing entry at
similarity \(\rho-\epsilon\), just outside the geometric deletion radius
\(\rho=0.85\). The baseline is geometric deletion alone, whose survival rate
is \pct{31.1}.

\begin{table}[H]
\centering
\caption{Synthetic verbatim row placed outside the geometric radius. Attack
gain is the fraction changed from wrong under baseline geometric deletion to
correct after insertion.}
\label{tab:adversarial}
\begin{tabular}{rrrr}
\toprule
\(\epsilon\) & Synthetic row selected & Post-attack correct & Attack gain \\
\midrule
0.01 & \pct{87.2} & \pct{59.3} & \pct{36.1} \\
0.02 & \pct{77.2} & \pct{58.3} & \pct{33.3} \\
0.05 & \pct{56.3} & \pct{50.9} & \pct{24.0} \\
\bottomrule
\end{tabular}
\end{table}

An explicit answer value immediately outside the radius is often selected and
restores the answer. The same qualitative effect appears in all five
datasets. This shows that the cosine threshold is a routing boundary, not a
semantic deletion guarantee. Hyphenated, letter-spaced, and prefix-cued
values add almost no new correct answers on T-REx (attack gain at most
\pct{0.2}). This is a synthetic verbatim boundary test, not evidence about
natural paraphrases or general obfuscation.

\subsection{Secondary diagnostics}

\paragraph{Representation probe.}
The probe is a five-fold closed-set ridge-ranking test, not an open-answer
decoder. It chooses among 315 candidate answers using event-0 \FULL\ query
vectors. It selects the accepted answer for 125 of 572 facts (\pct{21.9}),
compared with 41 of 572 (\pct{7.2}) for the most common answer. Behavioral
\DELOFF\ correctness is 130 of 572 (\pct{22.7}), but the two sets overlap
weakly (Jaccard \(\pct{19.2}\)). For 13 cohort cases, event 0 is not the later
audited answer-mention event, so this remains a descriptive diagnostic.

\paragraph{Radius sweep and collateral.}
On T-REx, geometric efficacy rises from \pct{61.7} at \(\rho=0.95\) to
\pct{76.8} at \(\rho=0.70\). Measured collateral is zero through
\(\rho=0.75\) and \pct{0.22} at \(\rho=0.70\). We do not treat this as a
headline result. In 450 of 499 target neighborhoods, the requested cosine
ball contained too few comparison facts and was filled with facts outside
the ball; only 552 of 2,607 evaluated links were inside it. The approximate
search-depth issue also applies to this sweep.

\section{What we think is worth testing next}

Let each memory entry be \(e_i=(k_i,v_i,s_i)\), containing a continuous key,
text value, and source metadata. The next useful baseline is a persistent
identity assignment \(z_i\) such that
\[
z_i=z_j
\quad\Longleftrightarrow\quad
e_i \text{ and } e_j \text{ express the same deletable proposition},
\]
under a clear deletion policy. A request for identity \(z\) should hide every
entry assigned to \(z\), while retaining entries that merely share an answer,
topic, or source.

For example, ``X was born in Paris'' and ``Y was born in Paris'' share an
answer but express different propositions. ``X's birthplace is Paris'' and
``X was born in the French capital'' may express the same proposition.
Identity therefore needs more than unconstrained clustering of key vectors.

\subsection{A minimal identity-aware baseline}

\begin{enumerate}
  \item Generate high-recall candidate pairs using key neighbors, entities,
  values, and source links.
  \item Label pairs as \emph{same proposition}, \emph{same answer only},
  \emph{related}, or \emph{unrelated}. Shared-answer examples are important
  hard negatives.
  \item Score proposition equivalence using text, entity and relation
  consistency, provenance, and key similarity. Form constrained clusters and
  store an inspectable identity ID.
  \item Delete by identity and rerun the same three-state audit with fixed
  search depth, cohort, and decoding.
\end{enumerate}

The key comparison is whether identity deletion approaches the oracle
filter's efficacy while preserving unrelated entries that share its answer.
We should report cluster precision and recall, cluster size, latency,
survival, retrieval-disabled answerability, retrieval interference, and
collateral on separate hard-negative sets for shared answer, entity,
relation, source, and key neighborhood. Splits should be identity-disjoint.

Calibrated abstention is also worth testing. The interference rates show that
accepting the next available candidate can turn a correct retrieval-disabled
answer into a wrong answer. Abstention and identity resolution address
different problems and can be evaluated separately.

\section{What this note does not claim}

\begin{itemize}
  \item The five main results are conditional on \FULL-correct cases with a
  pre-answer answer-mention retrieval. They do not measure all knowledge in
  the model or index.
  \item The answer-mention heuristic does not verify the full
  subject--relation--object proposition.
  \item Gold-answer filtering uses the target answer and aliases at inference
  time. It is an efficacy upper baseline, not a deployable deletion policy.
  \item \DELOFF\ removes factual retrieval but does not separate model
  parameters from information in the prompt.
  \item Search-time filtering tests behavioral deletion, not physical
  erasure, replicas, caches, or later index rebuilds.
  \item Approximate retrieval and policy-dependent search depth limit causal
  comparisons between deletion sets.
  \item These 1,500-example slices and our decoding and scoring rules are not
  direct replications of the paper's benchmark or TOFU experiments.
  \item Identity clustering, modified pretraining, and learned abstention have
  not yet been implemented.
\end{itemize}

\section{Our overall read}

The results are positive for the \colmlm\ direction: the external memory is
editable, inspectable, and clearly affects behavior. The open problem is the
unit of deletion. Selected-row and cosine policies miss many alternative
answer-bearing spans. Gold-answer filtering closes most observed paths, but
it is an oracle string rule and can be too broad. The audit uses event-aligned
retrieval traces, paired interventions, controls, five-dataset coverage, and
concrete baselines for the next step: proposition-level identity keys with
measured efficacy and collateral damage.

\medskip
\noindent\textbf{Statistical note.} Intervals are two-sided 95\% Wilson
intervals. Paired comparisons use two-sided exact McNemar tests; exploratory
tests are not adjusted for multiplicity.

\begin{thebibliography}{9}

\bibitem{feldman2026colmlm}
Yair Feldman, Linxi Zhao, Nathan Godey, Dongyoung Go, Yilun Hua,
Kilian Q. Weinberger, Jennifer J. Sun, and Yoav Artzi.
\newblock Co-LMLM: Continuous-Query Limited Memory Language Models.
\newblock \emph{arXiv preprint arXiv:2607.07707}, 2026.
\newblock \url{https://arxiv.org/abs/2607.07707}.

\bibitem{maini2024tofu}
Pratyush Maini, Zhili Feng, Avi Schwarzschild, Zachary C. Lipton, and
J. Zico Kolter.
\newblock TOFU: A Task of Fictitious Unlearning for LLMs.
\newblock \emph{arXiv preprint arXiv:2401.06121}, 2024.
\newblock \url{https://arxiv.org/abs/2401.06121}.

\bibitem{raeesi2026audit}
Arya Raeesi and Hanna Roed.
\newblock Auditing Forgetting in Limited Memory Language Models.
\newblock \emph{arXiv preprint arXiv:2607.00605}, 2026.
\newblock \url{https://arxiv.org/abs/2607.00605}.

\bibitem{tavakoli2026deletion}
Leila Tavakoli and Mark Sanderson.
\newblock Deletion Isn't Enough: Auditing RAG for Selective Forgetting.
\newblock In \emph{Proceedings of the 49th International ACM SIGIR Conference
on Research and Development in Information Retrieval (SIGIR '26)},
pages 2721--2733, 2026.
\newblock \url{https://doi.org/10.1145/3805712.3808545}.

\end{thebibliography}

\end{document}
